\hl{Bangla is spoken by more than 270 million people \cite{jawad2026fivedialects}, but most
language technology for Bangla understands only one form of it: Standard
Bangla. As with many languages spoken by many people, naturally, there
are many variations to how Bangla is spoken. These variations are
dialects and this creates a communication barrier between people from
different regions. At home and online, many people write the way they
speak, in a regional dialect such as Chittagonian or Sylheti, or type
Bangla with Latin letters (romanized Bangla). This chapter explains why
that matters for question answering (QA) in healthcare and education,
states the problem, the research questions and the objectives,
summarises the method and the contributions, and sets out the scope and
the main challenges of the work.}

\section{Background}\label{background}

\hl{Ranking among the world's most prominent languages by total
speakers, Bangla has a significant global presence; across Bangladesh,
India and the diaspora, millions of people use it as their first or
second language \cite{shafayat2024benqa}, \cite{jawad2026fivedialects}. Bangla has a strong regional dialect
continuum: Bangladesh alone has about 55 regional varieties spread over
its 64 districts \cite{faria2023vashantor}, \cite{samin2024bangladialecto}. These varieties differ from Standard
Bangla (cholito bhasha), the form used in textbooks, news and
administration, in pronunciation, vocabulary, verb morphology and word
order. Many speakers learn a regional variety at home and Standard
Bangla only at school. Chittagonian and Sylheti are two of the most
widely spoken and most divergent varieties: in dialect translation and
dialect identification experiments they are consistently among the
hardest to handle \cite{jawad2026fivedialects}, \cite{faria2023vashantor}, \cite{khandaker2024bridging}, \cite{prantik2025}. In our own data, a
Chittagonian sentence and its Standard translation share only 30\% of
their character trigrams, and a Sylheti sentence 35\% (Chapter 4).}

\hl{Resources for natural language processing (NLP) are unevenly
distributed across the world's languages, and low-resource languages
receive only a small share of data, models and benchmarks \cite{joshi2020state}.
Bangla has gained important resources in recent years: pretrained
encoders and sequence-to-sequence models such as BanglaBERT and BanglaT5
\cite{bhattacharjee2022banglabert}, \cite{bhattacharjee2023banglanlg}, Bangla large language models (LLMs) such as TigerLLM
and TituLLMs \cite{raihan2025tigerllm}, \cite{titulm2025}, and QA benchmarks such as BanglaRQA,
NCTB-QA, BEnQA and SOMAJGYAAN \cite{ekram2022banglarqa}, \cite{nctbqa}, \cite{shafayat2024benqa}, \cite{shifa2025somajgyaan}. All
of these, and the textbooks, news articles, encyclopedia passages and
exam papers they are built from, are in Standard Bangla, and models are
evaluated only on Standard Bangla \cite{jawad2026fivedialects}. Figure~\ref{fig:ch1-landscape} summarises this
situation.}

% ---------------- Figure 1.1 (Background) ----------------
\begin{figure}[!htbp]
\centering
\begin{tikzpicture}[node distance=3mm and 4mm,
  res/.style={font=\footnotesize, text width=28mm, minimum height=12mm, execute at begin node={\hyphenpenalty=10000}},
  hdr/.style={font=\footnotesize\bfseries, text=ink, anchor=west}]
\node[hdr] (h1) at (0,0) {(a) Bangla NLP resources: built and evaluated on Standard Bangla};
\node[box, res, below=2mm of h1.south west, anchor=north west] (r1a) {\textbf{Encoders}\\BanglaBERT,\\BanglaT5};
\node[box, res, right=of r1a] (r1b) {\textbf{Bangla LLMs}\\TigerLLM,\\TituLLMs};
\node[box, res, right=of r1b] (r1c) {\textbf{QA benchmarks}\\BanglaRQA, NCTB-QA, BEnQA};
\node[box, res, right=of r1c] (r1d) {\textbf{Knowledge}\\textbooks, health\\questions and answers};
\node[hdr, below=6mm of r1a.south west, anchor=north west] (h2) {(b) Dialect resources: parallel text for translation, no question answering};
\node[safe, res, below=2mm of h2.south west, anchor=north west] (r2a) {\textbf{Parallel corpora}\\Vashantor, ONUBAD, Kothon, ChatgaiyyaAlap};
\node[safe, res, right=of r2a] (r2b) {\textbf{Dialect health questions}\\BanglaCHQ-Prantik};
\node[safe, res, right=of r2b] (r2c) {\textbf{LLM dialect benchmarks}\\lowest scores for Chittagonian};
\node[grey, res, right=of r2c] (r2d) {\textbf{Dialect QA data}\\health or education:\\\textit{none available}};
\node[hdr, below=6mm of r2a.south west, anchor=north west] (h3) {(c) This thesis: dialect-aware, retrieval-grounded question answering};
\node[cpu, res, below=2mm of h3.south west, anchor=north west] (r3a) {\textbf{Question}\\Standard, Chittagonian, Sylheti or romanized};
\node[cpu, res, right=of r3a] (r3b) {\textbf{Route and normalize}\\dialect $\rightarrow$ Standard Bangla};
\node[cpu, res, right=of r3b] (r3c) {\textbf{Ground}\\domain knowledge base + LoRA expert};
\node[lora, res, right=of r3c] (r3d) {\textbf{Answer}\\Standard Bangla, evaluated per variety};
\foreach \a/\b in {r3a/r3b,r3b/r3c,r3c/r3d}{\draw[arr] (\a) -- (\b);}
\end{tikzpicture}
\caption[Bangla NLP resources, dialect resources and the gap addressed by this thesis.]{(a) Bangla language technology is built on Standard Bangla, while (b) dialect resources support translation but not question answering. (c) This thesis connects the two: dialect questions are normalized into Standard Bangla so that existing Standard-Bangla QA data and knowledge sources can be reused.}
\label{fig:ch1-landscape}
\end{figure}

\hl{Large language models have made open-ended question answering widely
accessible, but they can produce fluent answers that no evidence
supports, a behaviour known as hallucination \cite{ji2023survey}.
Retrieval-augmented generation (RAG) reduces this risk by retrieving
relevant documents at answer time and conditioning the generator on them
\cite{lewis2020rag}. Parameter-efficient fine-tuning with low-rank adaptation
(LoRA) and its quantized variant QLoRA makes it possible to adapt
billion-parameter models on a single GPU \cite{hu2022lora}, \cite{dettmers2023qlora}. The idea of
dividing a task among specialised sub-models chosen by a gating or
routing function goes back to the mixture-of-experts (MoE) literature
\cite{jacobs1991adaptive}, \cite{shazeer2017moe}. Together, these techniques make a modular, grounded
and affordable QA system realistic for a low-resource language.}

\hl{The dialect problem matters most where reliable information is
scarce. In \textbf{healthcare}, Bangladesh has about 9.9 doctors, nurses
and midwives per 10,000 people, and only about 11\% of doctors and
nurses work in rural areas, where about two thirds of the population
live \cite{nuruzzaman2022informing}. Digital health tools mostly serve European users and
languages, and a Bangladeshi triage application reported that a general
speech-to-text service struggled with dialect speech \cite{nahar2025amardoctor}. Real consumer
health questions in Bangla are informal, loosely described and
code-mixed with English medical terms \cite{khan2023nervous}, Bangla telemedicine
consultations are full of vague descriptions, delayed disclosure and
regional expressions, and prioritising patients from such real
consultations is hard even for strong models (best triage macro-F1
0.463) \cite{saha2026doctalkbn}. In \textbf{education}, university students in
Bangladesh already use LLMs for learning \cite{hossain2025ailiteracy}, and
national-curriculum (NCTB) textbooks have become the basis of new Bangla
QA resources and models \cite{raihan2025tigerllm}, \cite{nctbqa}. Yet LLMs answer the same
exam questions less accurately in Bangla than in English \cite{shafayat2024benqa},
earlier Bangla textbook QA systems were extractive models that only
worked when the correct passage was supplied \cite{khondoker2024unlocking}, and no
educational QA resource contains questions asked in a regional dialect.}

\section{Rationale of the Study}\label{rationale-of-the-study}

\hl{\textbf{The dialect gap is large and prompting alone does not close
it.} A benchmark of 19 LLMs on questions in nine Bengali dialects found
far lower scores for Chittagonian (5.44 out of 10) than for Tangail
(7.68) \cite{sami2026dialectbias}. When LLMs are prompted to convert dialect text into
Standard Bangla, some score below a baseline that simply copies the
input unchanged \cite{jawad2026fivedialects}. In the reverse direction (Standard Bangla to
dialect), the best prompted system on dialect health questions reached a
BLEU of only 23.67 for Sylheti, and Chittagonian scores were generally
lower (14.78 in the same setting) \cite{prantik2025}; without retrieved examples,
LLMs failed almost completely at standard-to-dialect translation (BLEU
5--12) \cite{sami2025comparative}. In our own paired evaluation, the same questions asked
in dialect scored 33\% (Chittagonian) and 43\% (Sylheti) lower than in
Standard Bangla when the system did not normalize them (Chapter 5). If a
system does not understand the question, the quality of its knowledge
base and generator hardly matters.}

\hl{\textbf{Small amounts of targeted supervision help.} LoRA
fine-tuning of open LLMs on only 160 examples per dialect raised
Mistral-7B from 16.2 to 73.6 BLEU on dialect-to-English translation,
above every prompted closed-source model \cite{jawad2026fivedialects}. A lightly adapted
dialect component is therefore a realistic way to close much of the gap
at low cost.}

\hl{\textbf{Normalizing into Standard Bangla reuses existing resources.}
Almost all Bangla QA data, textbooks and health information exist only
in Standard Bangla, and building separate datasets and knowledge bases
for every dialect is not realistic. Rewriting a dialect question into
Standard Bangla first lets the system reuse all of them. There is
supporting evidence for this ``translate into the better-resourced
variety'' strategy: appending an English translation of the question
improved LLM accuracy on BEnQA across all subjects \cite{shafayat2024benqa}, and
dialect-to-Standard translation is the central step of Bangla
speech-standardization pipelines \cite{samin2024bangladialecto}.}

\hl{\textbf{Different domains need different evidence and safety
behaviour.} Health answers must be grounded in reviewed information, and
the system must recognise emergencies and know when to refer the user to
a doctor. Education answers should follow the curriculum, and the system
must recognise questions that its sources cannot answer; NCTB-QA makes
42.75\% of its questions unanswerable for this reason \cite{nctbqa}. Routing
questions to domain-specific pipelines keeps the prompts, knowledge
bases and safety rules of each domain separate and easy to inspect, and
small routers that choose a retrieval strategy per question are an
established idea \cite{jeong2024adaptiverag}. Showing users the passages behind each answer
is a simple way to make the system's reasoning inspectable.}

\hl{\textbf{The system must be affordable.} The institutions that would
benefit most, such as rural clinics, schools and small organisations,
cannot run very large models or pay per-question API fees. A design with
CPU routers, one shared open 4-billion-parameter model and a few small
LoRA adapters can be trained and served on one rented GPU \cite{dettmers2023qlora}. This
approach keeps the system easy enough to build for a student project,
yet lightweight enough for constrained budgets.}

\section{Problem Statement}\label{problem-statement}

\hl{Bangla QA systems and LLMs are developed and evaluated mainly on
Standard Bangla \cite{raihan2025tigerllm}, \cite{ekram2022banglarqa}, \cite{nctbqa}, \cite{shafayat2024benqa}. When a question
is written in Chittagonian or Sylheti, or in romanized Bangla, three
problems arise. First, dialect words, verb forms and spellings may be
misunderstood, and because the question shares few words with Standard
Bangla documents, retrieval may fail as well \cite{jawad2026fivedialects}, \cite{prantik2025}. Second,
without grounding in domain-appropriate evidence, a general model may
answer health and education questions from memory, which risks
hallucinated or unsafe answers \cite{ji2023survey}, \cite{saha2026doctalkbn}. Third, the
performance gap between Standard Bangla users and dialect users is not
measured, because existing benchmarks do not report results by language
variety.}

\hl{Formally, a user question \emph{q} is written in a variety \emph{v}
\ensuremath{\in} \{Standard, Chittagonian, Sylheti\} (in Bangla or Latin script) and
belongs to a domain \emph{d} \ensuremath{\in} \{education, general knowledge, medical,
out-of-scope\}. The system must identify \emph{v} and \emph{d}, produce
a Standard Bangla version of \emph{q} when \emph{v} is a dialect or the
text is romanized, retrieve evidence from the knowledge base of
\emph{d}, and return a Standard Bangla answer grounded in that evidence;
or it must ask for clarification, abstain, or give an emergency referral
when an answer would not be reliable or safe. Answer quality should be
as similar as possible across the three varieties.}

\hl{This thesis answers the following research questions:}

\begin{quote}
\hl{\textbf{RQ1} How accurately, and with how well-calibrated
confidence, can a lightweight router identify the language variety and
the domain of a Bangla question?}

\hl{\textbf{RQ2} How well can a parameter-efficiently fine-tuned model
normalize Chittagonian, Sylheti and romanized questions into Standard
Bangla, and how much of the dialect gap in retrieval and answer quality
does this normalization close?}

\hl{\textbf{RQ3} How much do retrieval, fine-tuned domain experts,
domain routing and the answer/clarify/abstain controller contribute to
answer quality and behaviour, and do routed experts outperform a single
shared adapter?}

\hl{\textbf{RQ4} Which open base model offers the best balance of answer
quality, dialect handling and compute cost under an identical
fine-tuning budget, and does Bangla-specific pre-training help after
fine-tuning?}
\end{quote}

\section{Objectives}\label{objectives}

\hl{The main objective is to design, implement and evaluate a
dialect-aware, multi-domain Bangla QA framework that accepts questions
in Standard Bangla, Chittagonian or Sylheti, in Bangla or Latin script,
and returns grounded Standard Bangla answers for education,
general-knowledge and medical questions. The main deliverables are
trained routing models, a dialect normalization expert, domain-specific
RAG experts, an evaluation suite that reports every result separately
for each language variety, and a reproducible code base. The specific
objectives are:}

\begin{itemize}
\item \hl{to build lightweight, calibrated routers that identify the
language variety and the domain of a question \cite{faria2023vashantor}, \cite{bhattacharjee2022banglabert};}

\item \hl{to train a dialect normalization expert that rewrites
Chittagonian, Sylheti and romanized questions into Standard Bangla with
parameter-efficient fine-tuning \cite{jawad2026fivedialects}, \cite{hu2022lora};}

\item \hl{to build a separate knowledge base and retriever for each domain,
and to adapt one shared open base model to each domain with LoRA/QLoRA
experts \cite{lewis2020rag}, \cite{dettmers2023qlora};}

\item \hl{to decide, for every question, whether to answer, ask for
clarification or abstain, and to handle medical emergencies and
medicine-dose requests with deterministic safety rules and a
doctor-referral note \cite{rajpurkar2018squad2}, \cite{kamath2020selective};}

\item \hl{to compare seven open base models, including Gemma 4 and two
Bangla-adapted models (TigerLLM and TituLM) with their original
versions, under an identical fine-tuning budget \cite{raihan2025tigerllm}, \cite{titulm2025},
\cite{gemma4};}

\item \hl{to evaluate every component and the complete system, using
NCTB-QA, BanglaRQA and dialect resources, with results reported
separately for Standard Bangla, Chittagonian and Sylheti, with
ablations, confidence intervals and significance tests, and to test the
system on dialect questions written by native speakers \cite{nctbqa},
\cite{ekram2022banglarqa}.}
\end{itemize}

\section{Methodology in Brief}\label{methodology-in-brief}

\hl{The framework follows a \emph{route, normalize, then ground} design
(Figure~\ref{fig:pipeline} in Chapter 4). A question first passes rule-based safety
checks for emergencies and dose requests. A calibrated linear classifier
on character and word n-grams identifies the variety; if the question is
Chittagonian or Sylheti with enough confidence, or romanized, a
normalization expert rewrites it into Standard Bangla. A second
classifier chooses the domain, and a controller decides whether to
answer, ask for clarification or abstain. For answerable questions, the
retriever of the chosen domain searches that domain's knowledge base
(BM25 for education and general knowledge, dense BGE-M3 retrieval for
medical questions \cite{robertson2009bm25}, \cite{chen2024bgem3}), and the domain's LoRA expert
writes the answer from the three best passages. All experts are small
adapters on one 4-bit Gemma-4-E4B-it model \cite{gemma4}, a
\emph{Mixture-of-Adapters} \cite{hu2022lora}, \cite{houlsby2019adapters}, \cite{shazeer2017moe}. An output guard
adds a doctor-referral note to medical answers and refuses to give
medicine doses. Figure~\ref{fig:ch1-overview} gives an overview of
this pipeline.}

% ---------------- Figure 1.2 (Methodology in Brief) ----------------
\begin{figure}[!htbp]
\centering
\begin{tikzpicture}[node distance=9mm and 6mm,
  st/.style={font=\footnotesize, text width=25.5mm, minimum height=14mm, execute at begin node={\hyphenpenalty=10000}}]
\node[grey, st] (q) {\textbf{Question}\\any variety, Bangla or Latin script};
\node[safe, st, right=of q] (s) {\textbf{Safety check}\\emergency, dose request};
\node[cpu, st, right=of s] (vr) {\textbf{Variety router}\\Standard / Chittagonian / Sylheti};
\node[lora, st, right=of vr] (n) {\textbf{Normalization expert}\\dialect $\rightarrow$ Standard};
\node[cpu, st, right=of n] (dr) {\textbf{Domain router + controller}\\answer, clarify or abstain};
\draw[arr] (q) -- (s); \draw[arr] (s) -- (vr);
\draw[arr] (vr) -- node[lab, above=1pt, font=\scriptsize]{dialect} (n); \draw[arr] (n) -- (dr);
\draw[darr] (vr.north) to[out=40,in=140] node[lab, above, font=\scriptsize]{Standard: no rewriting} (dr.north);
\node[box, st, below=of dr] (r) {\textbf{Domain retrieval}\\BM25 or BGE-M3, top 3 passages};
\node[lora, st, below=of n] (e) {\textbf{Routed LoRA expert}\\on one 4-bit Gemma-4-E4B};
\node[safe, st, below=of vr] (g) {\textbf{Output guard}\\referral note, no doses};
\node[grey, st, below=of s] (a) {\textbf{Answer}\\Standard Bangla + evidence};
\draw[arr] (dr) -- node[lab, right, font=\scriptsize]{answer} (r);
\draw[arr] (r) -- (e); \draw[arr] (e) -- (g); \draw[arr] (g) -- (a);
\draw[darr] (s) -- node[lab, left, font=\scriptsize, align=right]{emergency:\\referral reply} (a);
\end{tikzpicture}
\caption[Overview of the proposed route, normalize, then ground framework.]{Overview of the proposed \emph{route, normalize, then ground} framework. The routers and the controller run on the CPU; the normalization expert and the domain experts are LoRA adapters on one shared 4-bit model (a Mixture-of-Adapters). Each component is described in Chapter 4 (Figure~\ref{fig:pipeline}).}
\label{fig:ch1-overview}
\end{figure}

\hl{The work was carried out in stages. The data stage cleaned and split
more than twenty public data sources, including four dialect corpora
(Vashantor, Kothon, ChatgaiyyaAlap and ONUBAD), two QA datasets with
unanswerable questions (NCTB-QA and BanglaRQA), textbooks and medical
resources \cite{faria2023vashantor}, \cite{kothon2026}, \cite{chatgaiyyaalap2025}, \cite{sultana2025onubad}, \cite{nctbqa}, \cite{ekram2022banglarqa},
\cite{khan2023nervous}, \cite{aurpa2026banglameder}. Seven open models, including the Bangla-adapted
TigerLLM and TituLM and the multilingual Gemma-4-E4B, Qwen3 and
Llama-3.2, were then compared under an identical QLoRA budget \cite{raihan2025tigerllm},
\cite{titulm2025}, \cite{gemma4}, \cite{qwen3}, \cite{meta2024llama32}, \cite{dettmers2023qlora}. The experts, routers
and retrievers were trained and chosen on validation data, and the test
data were used only once, for reporting. Finally, the complete system
was evaluated under nine conditions on 2,225 paired test questions (780
Standard questions with 720 Chittagonian and 725 Sylheti versions), and
on 481 dialect questions written by native speakers. Every result is
reported separately for each language variety.}

\section{Contributions}\label{contributions}

\begin{enumerate}
\item \hl{\textbf{A modular dialect-aware QA framework} for Bangla that
separates variety routing, dialect normalization, domain routing,
answer/clarify/abstain control, per-domain retrieval, routed LoRA
experts on one shared base model, and medical safety rules, together
with its complete, reproducible implementation.}

\item \hl{\textbf{A fine-tuned dialect normalization expert} for
Chittagonian, Sylheti and romanized input that raises chrF++ from 51.2
to 76.3 (Chittagonian) and from 55.4 to 74.0 (Sylheti) after one epoch
(a second epoch adds 1.3 chrF++ overall), transfers to three dialects
never seen in training, beats similarity-based few-shot prompting, and
on the official Vashantor test sentences equals or exceeds published
dedicated translation models except one (Chittagonian BLEU against
DialectBanglaT5).}

\item \hl{\textbf{A paired, per-variety evaluation protocol} (the same
questions in three varieties, nine system conditions, Holm-corrected
significance tests) that measures the dialect gap end to end. It shows
that normalization closes 61\% (Chittagonian) and 64\% (Sylheti) of the
gap without harming Standard Bangla input; to our knowledge, no earlier
Bangla work measured this downstream effect.}

\item \hl{\textbf{An evaluation on 481 dialect questions written by native
speakers}, which shows that with normalization the system treats a
dialect question like its Standard version 86.1\% of the time (65.3\%
without).}

\item \hl{\textbf{A controlled model comparison and informative negative
results}: under an identical budget, Gemma-4-E4B-it is the best of seven
open models; Bangla continued pre-training gives no advantage after
fine-tuning; routed experts equal a single pooled adapter in accuracy,
so their value is modularity; and a contrastively trained
dialect-invariant retriever is not better overall than its base encoder
and does not reduce the dialect gap.}
\end{enumerate}

\section{Scope and Challenges}\label{scope-and-challenges}

\hl{This thesis builds and evaluates a research prototype, not a
ready-to-use medical product; both its individual parts and the complete
system are evaluated. Its scope and main challenges are listed below.}

\textbf{\hl{Scope of the Study}}

\begin{itemize}
\item \hl{\textbf{Input:} text questions in Standard Bangla, Chittagonian or
Sylheti, written in Bangla script or romanized. Three further dialects
(Barishal, Mymensingh, Noakhali) are used only as unseen test data.
Speech input is not part of this study, although a Bangla medical speech
recogniser with Sylheti recordings \cite{kabir2024asr} could be placed in front of
the system later.}

\item \hl{\textbf{Domains:} education (NCTB curriculum), general knowledge
(Bangla Wikipedia passages of BanglaRQA) and medical questions, each
with its own knowledge base and retrieval pipeline; other questions are
recognised and declined.}

\item \hl{\textbf{Output:} answers in Standard Bangla, with the evidence
shown to the user, because the knowledge sources and most readers use
the standard form. Answering back in the user's dialect is left for
future work.}

\item \hl{\textbf{Healthcare:} general health information only, with
emergency referral and a doctor-referral note; the system does not
diagnose or prescribe.}

\item \hl{\textbf{Compute:} open models trained with QLoRA on one rented
NVIDIA A40 GPU, with all non-GPU work on free Google Colab CPU runtimes
\cite{hu2022lora}, \cite{dettmers2023qlora}.}
\end{itemize}

\textbf{\hl{Key Challenges}}

\begin{itemize}
\item \hl{\textbf{Scarce and uneven dialect data:} Parallel dialect data are
limited and uneven across dialects, and most of them are everyday
conversation rather than health or school questions \cite{faria2023vashantor}, \cite{prantik2025},
\cite{kothon2026}, so domain vocabulary such as symptom names is
under-represented. For example, the medical dialect extension of
BanglaCHQ-Summ contains 2,350 Sylheti but only 500 Chittagonian
questions \cite{prantik2025}.}

\item \hl{\textbf{No standard dialect spelling:} Chittagonian and Sylheti
have no agreed written form, so the same word appears with several
spellings, and romanized text shares almost no words with Bangla-script
text \cite{jawad2026fivedialects}.}

\item \hl{\textbf{No dialect QA test data:} No existing dataset contains
health or education questions in Chittagonian or Sylheti with reference
answers, so the dialect gap has to be measured with generated paired
questions and with human-written dialect questions without reference
answers.}

\item \hl{\textbf{Error propagation:} A wrong variety or domain decision, or
a poor normalization, affects every later stage; the system therefore
uses calibrated confidences and a controller rather than trusting each
component blindly.}

\item \hl{\textbf{Medical safety:} Wrong, overconfident or delayed health
answers can cause harm, and triage is hard even for large models
\cite{saha2026doctalkbn}; emergencies and dose requests must be handled
deterministically.}

\item \hl{\textbf{Unanswerable and vague questions:} Many questions cannot
be answered from the knowledge base or are under-specified; the system
must abstain or ask instead of guessing \cite{rajpurkar2018squad2}, \cite{kamath2020selective}, \cite{nctbqa}.}

\item \hl{\textbf{Evaluating free-text answers:} Medical answers are free
text, and word-overlap metrics under-rate correct answers that use
different words; several metrics, a check of grounding and human review
are needed.}

\item \hl{\textbf{Limited compute and budget:} All training had to fit one
rented GPU and a small budget (about 40 GPU-hours, roughly US\$20),
which makes quantized, parameter-efficient training essential \cite{dettmers2023qlora}.}
\end{itemize}
